% Luc Destombe --- licence CC BY-NC-SA 4.0
% https://creativecommons.org/licenses/by-nc-sa/4.0/deed.fr
% Fichier autonome : compiler avec pdflatex, deux fois (signets, nombre de pages).
\documentclass[11pt,a4paper]{article}
\makeatletter
\newif\iflycee@unecolonne
\lycee@unecolonnetrue
\newif\iflycee@cm
\newif\iflycee@palatino

\RequirePackage[utf8]{inputenc}
\RequirePackage[T1]{fontenc}
\RequirePackage[french]{babel}
\RequirePackage{etoolbox}

\newcommand{\ShorthandsOff}{\@ifundefined{shorthandoff}{}{\shorthandoff{;:!?}}}
\newcommand{\ShorthandsOn}{\@ifundefined{shorthandon}{}{\shorthandon{;:!?}}}
\ShorthandsOff
\AtBeginDocument{\ShorthandsOn}
\AtBeginEnvironment{tikzpicture}{\ShorthandsOff}

\RequirePackage{geometry}
\RequirePackage{fancyhdr}
\RequirePackage{lastpage}
\RequirePackage{multicol}
\RequirePackage{array}
\RequirePackage{enumitem}
\RequirePackage{xcolor}
\RequirePackage{needspace}

\RequirePackage{amsmath,amssymb}
\iflycee@cm\else\RequirePackage{fourier}\fi
\DeclareMathAlphabet{\mathcal}{OMS}{cmsy}{m}{n}
\renewcommand{\ttdefault}{lmtt}
\begingroup\catcode`\_=\active \catcode`\^=\active
\global\let_\sb \global\let^\sp
\endgroup
\AtBeginDocument{\catcode`\_=12 \mathcode`\_="8000
  \catcode`\^=12 \mathcode`\^="8000 }
\newlength{\retraitcalculs}
\setlength{\retraitcalculs}{2em}

\RequirePackage{tikz}
\usetikzlibrary{arrows.meta,calc,babel,shapes.misc}
\RequirePackage[most]{tcolorbox}
\tcbuselibrary{skins,breakable}

\definecolor{coulDef}{HTML}{1F4E79}
\definecolor{coulProp}{HTML}{7030A0}
\definecolor{coulRem}{HTML}{C55A11}
\definecolor{coulEx}{HTML}{2E7D32}
\definecolor{coulExo}{HTML}{B03A2E}
\definecolor{coulPart}{HTML}{1F4E79}
\definecolor{coulMeth}{HTML}{1B6E4F}
\definecolor{coulCourbe}{HTML}{0B5ED7}
\definecolor{coulCourbeB}{HTML}{C00000}
\definecolor{coulCourbeC}{HTML}{2E7D32}
\definecolor{coulGrille}{HTML}{8C8C8C}
\definecolor{coulRep}{HTML}{1B6E4F}
\definecolor{coulBar}{HTML}{2E7D32}

\newcommand{\R}{\mathbb{R}}
\newcommand{\Rs}{\mathbb{R}^{*}}
\newcommand{\N}{\mathbb{N}}
\newcommand{\Ns}{\mathbb{N}^{*}}
\newcommand{\Z}{\mathbb{Z}}
\newcommand{\Q}{\mathbb{Q}}
\renewcommand{\le}{\leqslant}
\renewcommand{\ge}{\geqslant}
\renewcommand{\approx}{\simeq}
\newcommand{\vect}[1]{\overrightarrow{#1}}
\newcommand{\norme}[1]{\left\lVert #1 \right\rVert}
\newcommand{\intervalle}[2]{\left[#1\,;#2\right]}
\RequirePackage{cancel}
\renewcommand{\CancelColor}{\color{coulExo}}
\newcommand{\fois}[1]{\times{\color{coulExo}#1}}
\newcommand{\barre}[1]{\cancel{#1}}
\newlength{\lycee@eqh}\newlength{\lycee@eqd}
\newsavebox{\lycee@eqa}\newsavebox{\lycee@eqb}\newsavebox{\lycee@eqc}
\newcommand{\lycee@eqcol}[1]{\hbox{$\renewcommand{\arraystretch}{1.35}%
  \begin{array}[t]{@{}>{\displaystyle}r@{}>{\displaystyle{}}l@{}}#1\end{array}$}}
\newcommand{\lycee@eqmesure}[1]{%
  \ifdim\ht#1>\lycee@eqh\setlength{\lycee@eqh}{\ht#1}\fi
  \ifdim\dp#1>\lycee@eqd\setlength{\lycee@eqd}{\dp#1}\fi}
\newcommand{\lycee@eqtrait}{\hspace{1.2em}\tikz[baseline=0pt]\draw[dash pattern=on 2pt off 2pt,
  line width=0.5pt] (0,\lycee@eqh) -- (0,-\lycee@eqd);\hspace{1.2em}}
\newcommand{\eqproduit}[2]{\par\smallskip\noindent C'est une équation produit :
  $A\times B=0 \iff A=0$ ou $B=0$.\par\nopagebreak\smallskip
  \begingroup\setlength{\lycee@eqh}{0pt}\setlength{\lycee@eqd}{0pt}%
  \setbox\lycee@eqa=\lycee@eqcol{#1}\setbox\lycee@eqb=\lycee@eqcol{#2}%
  \lycee@eqmesure\lycee@eqa\lycee@eqmesure\lycee@eqb%
  \noindent\hspace*{\retraitcalculs}\usebox\lycee@eqa\lycee@eqtrait\usebox\lycee@eqb\endgroup\par\smallskip}
\newcommand{\eqproduittrois}[3]{\par\smallskip\noindent C'est une équation produit :
  $A\times B\times C=0 \iff A=0$ ou $B=0$ ou $C=0$.\par\nopagebreak\smallskip
  \begingroup\setlength{\lycee@eqh}{0pt}\setlength{\lycee@eqd}{0pt}%
  \setbox\lycee@eqa=\lycee@eqcol{#1}\setbox\lycee@eqb=\lycee@eqcol{#2}\setbox\lycee@eqc=\lycee@eqcol{#3}%
  \lycee@eqmesure\lycee@eqa\lycee@eqmesure\lycee@eqb\lycee@eqmesure\lycee@eqc%
  \noindent\hspace*{\retraitcalculs}\usebox\lycee@eqa\lycee@eqtrait\usebox\lycee@eqb\lycee@eqtrait\usebox\lycee@eqc\endgroup
  \par\smallskip}

\tikzset{grille/.style={coulGrille,line width=0.4pt}}
\newcommand{\quadrillage}[4]{\draw[grille,step=1] (#1,#2) grid (#3,#4);}
\tikzset{
  croix/.style={cross out,draw,line width=0.6pt,inner sep=0pt,minimum size=4pt},
  croixdroite/.style={croix,rotate=45,line width=0.8pt},
}
\newcommand{\pt}[3]{\path (#1) node[croix]{} node[#2,font=\small,inner sep=2.5pt]{$#3$};}

\newcommand{\lycee@repere}[4]{%
  \draw[grille,step=1] (#1,#2) grid (#3,#4);
  \draw[-{Stealth[length=2mm]},thick] (#1,0)--({#3+0.4},0) node[below left=-1pt]{$x$};
  \draw[-{Stealth[length=2mm]},thick] (0,#2)--(0,{#4+0.4}) node[below left=-1pt]{$y$};
  \node[below left,font=\scriptsize,inner sep=1.5pt] at (0,0) {$0$};}
\newcommand{\lycee@gradx}[1]{\foreach \k/\l in {#1}{%
  \draw (\k,0.07)--(\k,-0.07) node[below,font=\scriptsize,inner sep=1.5pt]{$\l$};}}
\newcommand{\lycee@grady}[1]{\foreach \k/\l in {#1}{%
  \draw (0.07,\k)--(-0.07,\k) node[left,font=\scriptsize,inner sep=1.5pt]{$\l$};}}
\newcommand{\lycee@intff}[2]{\left[#1\,;#2\right]}
\newcommand{\lycee@intfo}[2]{\left[#1\,;#2\right[}
\newcommand{\lycee@intof}[2]{\left]#1\,;#2\right]}
\newcommand{\lycee@intoo}[2]{\left]#1\,;#2\right[}
\newcommand{\lycee@ens}[1]{\left\{#1\right\}}
\AtBeginDocument{%
  \providecommand{\repere}{\lycee@repere}%
  \providecommand{\gradx}{\lycee@gradx}%
  \providecommand{\grady}{\lycee@grady}%
  \providecommand{\Cf}{\mathcal{C}\sb{f}}%
  \providecommand{\Cg}{\mathcal{C}\sb{g}}%
  \providecommand{\intff}{\lycee@intff}%
  \providecommand{\intfo}{\lycee@intfo}%
  \providecommand{\intof}{\lycee@intof}%
  \providecommand{\intoo}{\lycee@intoo}%
  \providecommand{\ens}{\lycee@ens}}

\newlist{questions}{enumerate}{3}
\setlist[questions,1]{label=\arabic*\textdegree),leftmargin=*,itemsep=2pt,topsep=3pt}
\setlist[questions,2]{label=\alph*),leftmargin=*,itemsep=1pt,topsep=2pt}
\setlist[questions,3]{label=\roman*),leftmargin=*,itemsep=1pt,topsep=1pt}
\setlist[itemize]{label=\textbullet,leftmargin=*,itemsep=2pt,topsep=3pt}

\newcommand{\besoinplace}[1]{\par\penalty\@M\begingroup
  \dimen@=#1\relax\dimen@ii=\pagegoal\advance\dimen@ii-\pagetotal
  \ifdim\dimen@>\dimen@ii\ifdim\dimen@ii>\z@\vfil\fi\break\fi\endgroup}

\newcommand{\lycee@pied}{\thepage/\pageref{LastPage}}
\newcommand{\entete}[2]{%
  \pagestyle{fancy}\fancyhf{}%
  \fancyhead[L]{\footnotesize\color{coulPart}\textbf{#1}}%
  \fancyhead[R]{\footnotesize\color{coulPart}\textbf{#2}}%
  \fancyfoot[C]{\footnotesize\lycee@pied}%
  \renewcommand{\headrulewidth}{0.4pt}%
  \renewcommand{\headrule}{\hbox to\headwidth{%
    \color{coulPart!40}\leaders\hrule height \headrulewidth\hfill}}}

\RequirePackage{ccicons}
\newcommand{\lycee@licence}{{\scriptsize\color{gray}%
  \copyright\ Luc Destombe --- \ccbyncsa\ CC BY-NC-SA 4.0}}
\newif\iflycee@licence \lycee@licencetrue
\newcommand{\sanslicence}{\lycee@licencefalse}
\AtBeginDocument{\iflycee@licence\AddToHookNext{shipout/foreground}{%
  \put(0,0){\rlap{\kern\dimexpr1in+\hoffset+\oddsidemargin+\textwidth\relax
    \llap{\raisebox{-\dimexpr1in+\voffset+\topmargin+\headheight+\headsep
      +\textheight+\footskip\relax}{\lycee@licence}}}}}\fi}

\newcommand{\lycee@site}{faire-des-maths.fr}
\newcommand{\lycee@version}{{\scriptsize\color{gray}\lycee@site
  \ifdefined\lyceeversion\ --- Version~\lyceeversion\fi}}
\AtBeginDocument{\iflycee@licence\AddToHookNext{shipout/foreground}{%
  \put(0,0){\rlap{\kern\dimexpr1in+\hoffset+\oddsidemargin\relax
    \raisebox{-\dimexpr1in+\voffset+\topmargin+\headheight+\headsep
      +\textheight+\footskip\relax}{\lycee@version}}}}\fi}

\newcommand{\formulesserrees}{%
  \setlength{\abovedisplayskip}{3pt}\setlength{\belowdisplayskip}{3pt}%
  \setlength{\abovedisplayshortskip}{2pt}\setlength{\belowdisplayshortskip}{3pt}}

\setlength{\parindent}{0pt}

\RequirePackage[implicit=false,hidelinks,pdfencoding=auto,psdextra,bookmarksopen,
  bookmarksopenlevel=1,pdfcreator={},pdfproducer={}]{hyperref}
\RequirePackage{bookmark}
\hypersetup{pdfauthor={Luc Destombe},pdfkeywords={CC BY-NC-SA 4.0}}
\newcounter{lycee@signet}
\newcommand{\signet}[3][]{\stepcounter{lycee@signet}%
  \edef\lycee@dest{\if\relax\detokenize{#1}\relax
    signet.\arabic{lycee@signet}\else#1\fi}%
  \hypertarget{\lycee@dest}{}\bookmark[level=#2,dest=\lycee@dest]{#3}}
\pdfstringdefDisableCommands{%
  \def\R{R}\def\N{N}\def\Z{Z}\def\Q{Q}\def\vect#1{#1}\def\Cf{Cf}\def\Cg{Cg}\def\fois#1{×#1}%
  \def\barre#1{#1}\def\intervalle#1#2{[#1;#2]}\def\intff#1#2{[#1;#2]}%
  \def\textsuperscript#1{#1}\def\dfrac#1#2{#1/#2}\def\frac#1#2{#1/#2}%
  \def\vec#1{#1⃗}}%
\RequirePackage[official]{eurosym}

\tcbset{exercice corrige/.style={
  enhanced,breakable,before upper=\raggedright,
  colback=white,colframe=coulExo,
  boxrule=0.8pt,arc=2pt,
  left=6pt,right=6pt,top=8pt,bottom=5pt,
  before skip=9pt,after skip=4pt,
  coltitle=white,fonttitle=\bfseries\small,
  attach boxed title to top left={xshift=8pt,yshift=-\tcboxedtitleheight/2},
  boxed title style={colback=coulExo,arc=2pt,boxrule=0pt}}}
\newcounter{exo}
\newtcolorbox{exercice}[1][]{exercice corrige,
  phantom={\signet[exercice-\arabic{exo}]{2}{Exercice \arabic{exo}}},
  title={Exercice \theexo\ifx\relax#1\relax\else{} --- #1\fi}}
\newenvironment{exo}[1][\relax]{\refstepcounter{exo}\begin{exercice}[#1]}{\end{exercice}}

\geometry{top=1.8cm,bottom=1cm,left=1cm,right=1cm,headheight=14pt,footskip=0.6cm}
\renewcommand{\lycee@pied}{\thepage}

\renewcommand{\theexo}{\ifnum\value{exo}<10 0\fi\arabic{exo}}
\tcbset{exercice corrige/.style={
  enhanced,breakable,before upper=\raggedright,
  colback=white,colframe=coulExo,
  boxrule=0.9pt,arc=2pt,
  left=8pt,right=8pt,top=8pt,bottom=6pt,
  before skip=8pt,after skip=6pt,
  coltitle=white,fonttitle=\bfseries,
  attach boxed title to top left={xshift=8pt,yshift=-\tcboxedtitleheight/2},
  boxed title style={colback=coulExo,arc=2pt,boxrule=0pt}}}

\newcommand{\entetefiche}[2]{%
  \begin{center}
    \begin{tcolorbox}[enhanced,width=0.92\linewidth,colback=white,colframe=coulPart,
        boxrule=1.4pt,arc=3pt,halign=center,top=6pt,bottom=6pt]
      {\large\bfseries\color{coulPart} CORRIGÉ --- FICHE D'EXERCICES}\\[3pt]
      {\normalsize\bfseries #1}\\[2pt]
      {\normalsize #2}
    \end{tcolorbox}
  \end{center}}

\newcounter{partie}
\newcommand{\partie}[1]{%
  \stepcounter{partie}%
  \besoinplace{8\baselineskip}\vspace{4pt}%
  \begin{tcolorbox}[enhanced,colback=coulPart!8,colframe=coulPart,boxrule=0pt,
      leftrule=3pt,arc=0pt,left=6pt,right=4pt,top=3pt,bottom=3pt,
      before skip=6pt,after skip=2pt]
    \bfseries\color{coulPart}\Roman{partie} --- #1\signet{1}{\Roman{partie} --- #1}
  \end{tcolorbox}}
\newcommand{\partiesansnum}[1]{%
  \besoinplace{8\baselineskip}\vspace{4pt}%
  \begin{tcolorbox}[enhanced,colback=coulPart!8,colframe=coulPart,boxrule=0pt,
      leftrule=3pt,arc=0pt,left=6pt,right=4pt,top=3pt,bottom=3pt,
      before skip=6pt,after skip=2pt]
    \bfseries\color{coulPart}#1\signet{1}{#1}
  \end{tcolorbox}}

\setlist[questions,1]{label=\arabic*\textdegree),leftmargin=*,itemsep=3pt,topsep=3pt}
\setlist[questions,2]{label=\alph*),leftmargin=*,itemsep=2pt,topsep=2pt}
\setlist[questions,3]{label=\roman*),leftmargin=*,itemsep=1pt,topsep=1pt}
\newlist{expr}{enumerate}{1}
\setlist[expr]{label={\Alph*\ $=$},leftmargin=*,itemsep=2pt,topsep=3pt}

\newcommand{\res}[1]{{\color{coulRep}#1}}
\newcommand{\rem}[1]{\par\smallskip{\small\itshape\color{black!60}$\rightarrow$ #1}}
\newenvironment{calculs}{\par\smallskip\noindent$\displaystyle\begin{aligned}[t]}%
  {\end{aligned}$\par\smallskip}
\newcommand{\justif}[1]{\quad\text{\small\itshape\color{black!60}#1}}
\newcommand{\cas}[1]{\par\smallskip\noindent\textbf{\color{coulExo}#1}\vadjust{\nobreak}}
\newcommand{\explique}[2]{\par\smallskip\noindent
  \begin{minipage}[t]{0.6\linewidth}\raggedright #1\end{minipage}\hfill
  \begin{minipage}[t]{0.37\linewidth}\raggedright\leavevmode
    {\small\itshape\color{black!60}#2}\end{minipage}\par}

\setlength{\parskip}{2pt}

\makeatother
\entete{Chapitre 12 --- Statistiques --- corrigé}{Seconde}
\setlength{\parindent}{0pt}
\setlength{\jot}{5pt}

\let\casligne\cas
\renewcommand{\cas}[1]{\casligne{#1}\par\nopagebreak}

\newcommand{\cotefig}[3][0.40]{\par\noindent
  \begin{minipage}[t]{\dimexpr0.97\linewidth-#1\linewidth\relax}\raggedright #2\end{minipage}\hfill
  \begin{minipage}[t]{#1\linewidth}\vspace{0pt}\centering #3\end{minipage}\par}

\tikzset{barre/.style={fill=coulCourbe!20,draw=coulCourbe,line width=0.8pt},
  secteur/.style={draw=black,line width=0.6pt}}

\begin{document}

\entetefiche{Chapitre 12 --- Statistiques}{Seconde}

\partie{Effectifs, fréquences, effectifs cumulés, classes}

\begin{exo}[Activités hors du lycée]
\cas{a)}
\explique{Population : les \res{$25$ élèves qui ont répondu}.}{l'élève qui ne se
prononce pas est retiré}
\explique{Caractère : le nombre d'activités, \res{quantitatif discret}.}{des nombres
entiers, isolés}
\cas{b)}
\explique{On compte chaque valeur : six fois $0$, dix fois $1$, six fois $2$ et trois
fois $3$ ; $6+10+6+3=\res{25}$.}{on retrouve les $25$ réponses}
\begin{calculs}
f &= \dfrac{6}{25} &&\justif{fréquence de la valeur $0$}\\
f &= 0{,}24 &&\justif{soit $24\,\%$}
\end{calculs}
\begin{center}
\renewcommand{\arraystretch}{1.5}
\begin{tabular}{|l|*{4}{c|}c|}
\hline
Nombre d'activités & $0$ & $1$ & $2$ & $3$ & Total\\\hline
Effectif & $6$ & $10$ & $6$ & $3$ & $25$\\\hline
Fréquence & $24\,\%$ & $40\,\%$ & $24\,\%$ & $12\,\%$ & $100\,\%$\\\hline
Effectif cumulé croissant & $6$ & $16$ & $22$ & $25$ & \\\hline
\end{tabular}
\end{center}
\cas{c)}
\explique{L'effectif cumulé de la valeur $1$ vaut $16$ : \res{$16$ élèves} sont inscrits
à au plus une activité.}{au plus une : zéro ou une activité}
\begin{calculs}
f &= \dfrac{16}{25} &&\justif{fréquence cumulée de la valeur $1$}\\
f &= \res{0{,}64} &&\justif{soit $64\,\%$ des élèves}
\end{calculs}
\cas{d)}
\begin{calculs}
\text{effectif} &= 25-16 &&\justif{au moins deux : tous sauf « au plus une »}\\
\text{effectif} &= 9 &&
\end{calculs}
\explique{$\dfrac{9}{25}=0{,}36$ : \res{$36\,\%$} des élèves sont inscrits à au moins
deux activités.}{on peut aussi ajouter $24\,\%+12\,\%$}
\end{exo}

\begin{exo}[Le bon vocabulaire]
\cas{a)}
\explique{La couleur préférée est \res{qualitative} ; le nombre de messages, quantitatif
discret ; le temps sur $100$ m, quantitatif continu.}{un temps peut prendre toutes les valeurs d'un
intervalle}
\explique{La pointure est quantitative discrète ; la masse d'un cartable, quantitative
continue ; le sport pratiqué, qualitatif.}{la pointure ne prend que des valeurs isolées ($38$,
$38{,}5$, $39$\ldots)}
\cas{b)}
\explique{\res{La population est l'ensemble des élèves de seconde du lycée.}}{les
individus étudiés}
\explique{Le caractère étudié est la durée de sommeil ; il est quantitatif
continu.}{une durée}
\end{exo}

\begin{exo}[Effectifs cumulés d'un quiz]
\cas{a)}
\explique{Les élèves qui ont au plus $5$ ont eu trois, quatre ou cinq, et $1+2+4=\res{7}$.}{somme
des effectifs des notes inférieures ou égales à $5$}
\cas{b)}
\explique{On ajoute chaque effectif au précédent : $7+5=12$, puis $12+8=20$\ldots{}
La dernière case vaut \res{$32$}, l'effectif total.}{vérification}
\begin{center}
\renewcommand{\arraystretch}{1.5}
\begin{tabular}{|l|*{8}{c|}}
\hline
Note & $3$ & $4$ & $5$ & $6$ & $7$ & $8$ & $9$ & $10$\\\hline
Effectif & $1$ & $2$ & $4$ & $5$ & $8$ & $6$ & $4$ & $2$\\\hline
Effectif cumulé & $1$ & $3$ & $7$ & $12$ & $20$ & $26$ & $30$ & $32$\\\hline
\end{tabular}
\end{center}
\cas{c)}
\explique{Au plus $7$ : \res{$20$ élèves} (effectif cumulé de $7$).}{lecture directe}
\begin{calculs}
\text{effectif} &= 32-20 &&\justif{au moins $8$ : tous sauf ceux qui ont au plus $7$}\\
\text{effectif} &= \res{12} &&
\end{calculs}
\cas{d)}
\begin{calculs}
f &= \dfrac{12}{32} &&\justif{effectif cumulé de $6$ sur l'effectif total}\\
f &= 0{,}375 &&
\end{calculs}
\explique{\res{$37{,}5\,\%$ des élèves ont eu au plus $6$.}}{une phrase pour conclure}
\end{exo}

\begin{exo}[Venir au lycée]
\cas{a)}
\begin{calculs}
f &= 100\,\%-(44\,\%+24\,\%+18\,\%+8\,\%) &&\justif{la somme des fréquences vaut $100\,\%$}\\
f &= 100\,\%-94\,\% &&\\
f &= \res{6\,\%} &&
\end{calculs}
\cas{b)}
\begin{calculs}
\text{bus} &= 0{,}44\times 1~250 &&\justif{$44\,\%$ de $1~250$ (chapitre 7)}\\
\text{bus} &= \res{550} &&
\end{calculs}
\explique{De même : voiture $0{,}24\times 1~250=300$ ; à pied $225$ ; vélo $100$ ;
train $75$.}{vérification : $550+300+225+100+75=1~250$}
\cas{c)}
\explique{$12\,\%$ des élèves font $156$ élèves : $0{,}12\times N=156$.}{on cherche
l'effectif total $N$}
\begin{calculs}
N &= \dfrac{156}{0{,}12} &&\justif{on divise par $0{,}12$}\\
N &= \res{1~300} &&\justif{$1~300$ élèves}
\end{calculs}
\end{exo}

\begin{exo}[Distance au lycée]
\cas{a)}
\explique{On range chaque distance dans sa classe : $2$ va dans $\intfo{2}{4}$, $4$ dans
$\intfo{4}{6}$, $6$ dans $\intfo{6}{8}$.}{la borne de droite est exclue}
\begin{center}
\renewcommand{\arraystretch}{1.5}
\begin{tabular}{|l|*{5}{c|}c|}
\hline
Distance (km) & $\intfo{0}{2}$ & $\intfo{2}{4}$ & $\intfo{4}{6}$ & $\intfo{6}{8}$
  & $\intfo{8}{10}$ & Total\\\hline
Effectif & $5$ & $6$ & $4$ & $3$ & $2$ & $20$\\\hline
Fréquence & $25\,\%$ & $30\,\%$ & $20\,\%$ & $15\,\%$ & $10\,\%$ & $100\,\%$\\\hline
\end{tabular}
\end{center}
\explique{Vérification : $5+6+4+3+2=\res{20}$.}{aucune distance oubliée}
\cas{b)}
\begin{calculs}
f &= \dfrac{5}{20} &&\justif{classe $\intfo{0}{2}$}\\
f &= \res{25\,\%} &&
\end{calculs}
\explique{Les autres fréquences sont dans le tableau.}{$\dfrac{6}{20}=30\,\%$\ldots}
\cas{c)}
\explique{Moins de $4$ km : classes $\intfo{0}{2}$ et $\intfo{2}{4}$, soit
$25\,\%+30\,\%=\res{55\,\%}$ des élèves.}{on ajoute les fréquences}
\cas{d)}
\explique{\res{Non} : le tableau dit seulement que $6$ élèves habitent entre $2$ et
$4$ km, pas combien sont à moins de $3$ km.}{regrouper en classes fait perdre le
détail des valeurs}
\end{exo}

\partie{Moyenne, moyenne pondérée}

\begin{exo}[Trois moyennes]
\cas{a)}
\begin{calculs}
\overline{x} &= \dfrac{8+11+6+9+12+7+10}{7} &&\justif{$7$ jours}\\
\overline{x} &= \dfrac{63}{7} &&\\
\overline{x} &= \res{9} &&\justif{$9~000$ pas par jour en moyenne}
\end{calculs}
\cas{b)}
\begin{calculs}
\overline{x} &= \dfrac{6\times 0+10\times 1+6\times 2+3\times 3}{25}
  &&\justif{chaque valeur compte autant de fois que son effectif}\\
\overline{x} &= \dfrac{0+10+12+9}{25} &&\\
\overline{x} &= \dfrac{31}{25} &&\\
\overline{x} &= \res{1{,}24} &&\justif{activité par élève en moyenne}
\end{calculs}
\cas{c)}
\begin{calculs}
\overline{x} &= 0{,}55\times 1+0{,}3\times 2+0{,}1\times 3+0{,}05\times 4
  &&\justif{fréquences en écriture décimale}\\
\overline{x} &= 0{,}55+0{,}6+0{,}3+0{,}2 &&\\
\overline{x} &= \res{1{,}65} &&\justif{personne par voiture en moyenne}
\end{calculs}
\end{exo}

\begin{exo}[Moyenne avec coefficients]
\cas{a)}
\begin{calculs}
\overline{x} &= \dfrac{2\times 12+1\times 8+3\times 15+2\times 10}{2+1+3+2}
  &&\justif{les coefficients jouent le rôle des effectifs}\\
\overline{x} &= \dfrac{24+8+45+20}{8} &&\\
\overline{x} &= \dfrac{97}{8} &&\\
\overline{x} &= \res{12{,}125} &&
\end{calculs}
\explique{Sans coefficient, on trouverait $\dfrac{45}{4}=11{,}25$ : le $15$, de
coefficient $3$, compte trois fois, ce qui tire la moyenne vers le haut.}{la meilleure
note a le plus grand coefficient}
\cas{b)}
\explique{On note $x$ la note cherchée ; la somme des coefficients passe à $10$.}{$8+2$}
\begin{calculs}
\dfrac{97+2x}{10} &= 12{,}5 &&\justif{la moyenne voulue}\\
97+2x &= 125 &&\justif{on multiplie par $10$}\\
2x &= 28 &&\justif{on retire $97$}\\
x &= \res{14} &&\justif{on divise par $2$}
\end{calculs}
\explique{Inès doit obtenir $14$.}{vérification : $\dfrac{97+28}{10}=12{,}5$}
\end{exo}

\begin{exo}[Moyenne de deux classes]
\cas{a)}
\begin{calculs}
\overline{x} &= \dfrac{32\times 11{,}5+28\times 13}{60}
  &&\justif{chaque moyenne pèse son effectif}\\
\overline{x} &= \dfrac{368+364}{60} &&\\
\overline{x} &= \dfrac{732}{60} &&\\
\overline{x} &= \res{12{,}2} &&
\end{calculs}
\cas{b)}
\explique{\res{Les deux classes n'ont pas le même effectif} : la seconde A, plus
nombreuse, pèse davantage, et la moyenne est plus proche de $11{,}5$ que de $13$.}{
$\dfrac{11{,}5+13}{2}$ ne vaut que pour deux groupes de même effectif}
\cas{c)}
\begin{calculs}
\overline{x} &= \dfrac{732+30\times 10}{90} &&\justif{$732$ : somme des notes des $60$ premiers}\\
\overline{x} &= \dfrac{1~032}{90} &&\\
\overline{x} &\approx \res{11{,}47} &&
\end{calculs}
\end{exo}

\begin{exo}[Trajets des salariés]
\cas{a)}
\explique{Centres : $\dfrac{0+15}{2}=\res{7{,}5}$, puis $22{,}5$, $37{,}5$ et
$52{,}5$.}{milieu de chaque classe}
\cas{b)}
\begin{calculs}
\overline{x} &\approx \dfrac{18\times 7{,}5+30\times 22{,}5+24\times 37{,}5+8\times 52{,}5}{80}
  &&\justif{chaque classe est remplacée par son centre}\\
\overline{x} &\approx \dfrac{135+675+900+420}{80} &&\\
\overline{x} &\approx \dfrac{2~130}{80} &&\\
\overline{x} &\approx \res{26{,}6} &&\justif{$26{,}625$ arrondi : environ $26{,}6$ min}
\end{calculs}
\cas{c)}
\explique{Moins de $30$ min : $18+30=48$ salariés, et $\dfrac{48}{80}=\res{0{,}6}$, soit
$60\,\%$ des salariés.}{partie I : fréquence}
\end{exo}

\begin{exo}[Linéarité de la moyenne]
\cas{a)}
\explique{Chaque prix est multiplié par $1{,}05$, donc la moyenne aussi :
$24\times 1{,}05=\res{25{,}2}$ €.}{coefficient multiplicateur (chapitre 7) et
linéarité}
\cas{b)}
\explique{On retire $2$ € à chaque prix, donc à la moyenne :
$25{,}2-2=\res{23{,}2}$ €.}{propriété 1}
\cas{c)}
\begin{calculs}
\overline{x}_F &= 1{,}8\times 15+32 &&\justif{on applique la même transformation à la moyenne}\\
\overline{x}_F &= 27+32 &&\\
\overline{x}_F &= \res{59} &&\justif{$59$ degrés Fahrenheit}
\end{calculs}
\end{exo}

\begin{exo}[Une note de plus]
\cas{a)}
\begin{calculs}
\text{somme} &= 4\times 11 &&\justif{moyenne $=$ somme $\div 4$}\\
\text{somme} &= \res{44} &&
\end{calculs}
\cas{b)}
\begin{calculs}
\overline{x} &= \dfrac{44+16}{5} &&\justif{cinq notes}\\
\overline{x} &= \dfrac{60}{5} &&\\
\overline{x} &= \res{12} &&
\end{calculs}
\cas{c)}
\begin{calculs}
\dfrac{60+x}{6} &= 12{,}5 &&\justif{six notes}\\
60+x &= 75 &&\justif{on multiplie par $6$}\\
x &= \res{15} &&\justif{on retire $60$}
\end{calculs}
\end{exo}

\partie{Étendue, médiane, quartiles}

\begin{exo}[Médiane et quartiles de deux listes]
\cas{a)}
\explique{Série $A$ rangée : $19$ ; $21$ ; $22$ ; $24$ ; $25$ ; $26$ ; $27$ ; $28$ ; $29$ ;
$31$ ; $33$.}{on range d'abord, $N=11$}
\explique{Étendue : $33-19=\res{14}$ ans. $N$ impair : la médiane est la valeur de rang
$\dfrac{11+1}{2}=6$, $Me=\res{26}$ ans.}{valeur du milieu}
\explique{$\dfrac{11}{4}=2{,}75$, arrondi à $3$, donc $Q_1=\res{22}$ ;
$\dfrac{3\times 11}{4}=8{,}25$, arrondi à $9$, donc $Q_3=\res{29}$.}{rangs arrondis à
l'entier supérieur}
\cas{b)}
\explique{Série $B$ rangée : $0$ ; $0$ ; $1$ ; $1$ ; $2$ ; $2$ ; $2$ ; $3$ ; $4$ ; $5$.}{$N=10$}
\explique{Étendue : $5-0=\res{5}$. $N$ pair : moyenne des rangs $5$ et $6$,
$Me=\dfrac{2+2}{2}=\res{2}$.}{$\dfrac{10}{2}=5$}
\explique{$\dfrac{10}{4}=2{,}5$, arrondi à $3$, donc $Q_1=\res{1}$ ; $\dfrac{30}{4}=7{,}5$,
arrondi à $8$, donc $Q_3=\res{3}$.}{rangs $3$ et $8$}
\cas{c)}
\explique{\res{Lors d'au moins trois quarts des matchs, l'équipe a marqué au plus
$3$ buts.}}{$Q_3$ : au moins $75\,\%$ des valeurs lui sont inférieures ou égales}
\end{exo}

\begin{exo}[Séances de cinéma]
\cas{a)}
\explique{Effectifs cumulés : $8$ ; $20$ ; $29$ ; $35$ ; $38$ ; \res{$40$}.}{chaque
effectif s'ajoute au précédent}
\cas{b)}
\explique{$N=40$ : rangs $20$ et $21$. La $20^{\text{e}}$ valeur est $1$ (rangs $9$ à
$20$), la $21^{\text{e}}$ est $2$, donc $Me=\dfrac{1+2}{2}=\res{1{,}5}$.}{médiane d'un
effectif pair}
\explique{$\dfrac{40}{4}=10$ : la $10^{\text{e}}$ valeur est $1$, $Q_1=\res{1}$ ;
$\dfrac{3\times 40}{4}=30$ : la $30^{\text{e}}$ valeur est $3$ (rangs $30$ à $35$),
$Q_3=\res{3}$.}{rangs lus dans les effectifs cumulés}
\cas{c)}
\explique{Écart interquartile : $Q_3-Q_1=3-1=\res{2}$.}{définition 7}
\cas{d)}
\begin{calculs}
\overline{x} &= \dfrac{8\times 0+12\times 1+9\times 2+6\times 3+3\times 4+2\times 5}{40}
  &&\justif{partie II}\\
\overline{x} &= \dfrac{0+12+18+18+12+10}{40} &&\\
\overline{x} &= \dfrac{70}{40} &&\\
\overline{x} &= \res{1{,}75} &&
\end{calculs}
\end{exo}

\begin{exo}[Le salaire du patron]
\cas{a)}
\begin{calculs}
\overline{x} &= \dfrac{1~800+1~900+2~000+2~100+2~100+2~200+2~300+2~400}{8} &&\\
\overline{x} &= \dfrac{16~800}{8} &&\\
\overline{x} &= \res{2~100} &&\justif{euros}
\end{calculs}
\explique{Série déjà rangée, $N=8$ : $Me=\dfrac{2~100+2~100}{2}=\res{2~100}$ € ;
étendue : $2~400-1~800=600$ €.}{rangs $4$ et $5$}
\cas{b)}
\begin{calculs}
\overline{x} &= \dfrac{16~800+9~600}{9} &&\justif{neuf salaires}\\
\overline{x} &= \dfrac{26~400}{9} &&\\
\overline{x} &\approx \res{2~933{,}33} &&
\end{calculs}
\explique{$N=9$ : la médiane est la $5^{\text{e}}$ valeur, toujours $\res{2~100}$ € ;
étendue : $9~600-1~800=7~800$ €.}{le patron s'ajoute à la fin de la série}
\cas{c)}
\explique{\res{La médiane} : elle n'a pas bougé, alors que la moyenne a augmenté de plus de
$800$ € à cause d'un seul salaire, bien plus élevé que les autres.}{la moyenne est
sensible aux valeurs extrêmes, la médiane peu}
\end{exo}

\begin{exo}[Âges d'une association]
\cas{a)}
\explique{Effectifs cumulés : $6$ ; $20$ ; $29$ ; $37$ ; $44$ ; \res{$50$}.}{chaque effectif
s'ajoute au précédent}
\cas{b)}
\explique{$\dfrac{50}{2}=25$ et $20<25\le 29$ : la classe médiane est
\res{$\intfo{30}{40}$}.}{première classe où l'effectif cumulé atteint $25$}
\cas{c)}
\explique{Il faut $25-20=5$ membres de cette classe, sur ses $9$.}{les $20$ premiers ont
moins de $30$ ans}
\begin{calculs}
Me &\approx 30+\dfrac{5}{9}\times 10 &&\justif{début de la classe $+$ part de son amplitude}\\
Me &\approx 30+\dfrac{50}{9} &&\\
Me &\approx \dfrac{320}{9} &&\justif{$30=\dfrac{270}{9}$}\\
Me &\approx \res{35{,}6} &&\justif{environ $35{,}6$ ans}
\end{calculs}
\cas{d)}
\begin{calculs}
\overline{x} &\approx \dfrac{6\times 15+14\times 25+9\times 35+8\times 45+7\times 55+6\times 65}{50}
  &&\justif{centres des classes}\\
\overline{x} &\approx \dfrac{90+350+315+360+385+390}{50} &&\\
\overline{x} &\approx \dfrac{1~890}{50} &&\\
\overline{x} &\approx \res{37{,}8} &&\justif{ans}
\end{calculs}
\end{exo}

\begin{exo}[Contrôle d'une chaîne de remplissage]
\cas{a)}
\explique{On retire $500$ : les valeurs deviennent $-3$ ; $-2$ ; $-1$ ; $0$ ; $1$ ; $2$ ;
$3$, avec les mêmes effectifs.}{des nombres plus petits, plus faciles à calculer}
\begin{calculs}
\overline{y} &= \dfrac{2\times(-3)+5\times(-2)+14\times(-1)+40\times 0+24\times 1+10\times 2+5\times 3}{100} &&\\
\overline{y} &= \dfrac{-6-10-14+0+24+20+15}{100} &&\\
\overline{y} &= \dfrac{29}{100} &&\\
\overline{y} &= 0{,}29 &&
\end{calculs}
\explique{On avait retiré $500$ à chaque masse, donc la moyenne des masses vaut
$500+0{,}29=\res{500{,}29}$ g.}{linéarité de la moyenne}
\cas{b)}
\explique{Effectifs cumulés : $2$ ; $7$ ; $21$ ; $61$ ; $85$ ; $95$ ; $100$.}{$N=100$}
\explique{Les valeurs de rangs $50$ et $51$ valent $500$ g, donc $Me=\res{500}$ ; celle
de rang $25$ aussi, donc $Q_1=\res{500}$ ; celle de rang $75$ vaut $501$ g, donc
$Q_3=\res{501}$.}{rangs $22$ à $61$ : masse $500$ g ; rangs $62$ à $85$ : masse $501$ g}
\cas{c)}
\explique{En dehors de $\intff{498}{502}$ : les pots de $497$ g et de $503$ g, soit
$2+5=7$ pots sur $100$, \res{$7\,\%$}.}{$498$ et $502$ sont dans l'intervalle}
\cas{d)}
\explique{La moyenne est à $0{,}29$ g de $500$ g (moins de $0{,}5$) et $7\,\%<10\,\%$ :
\res{la machine est bien réglée}.}{les deux conditions sont remplies}
\end{exo}

\partie{Écart type}

\begin{exo}[Lectures régulières]
\cas{a)}
\begin{calculs}
\overline{x} &= \dfrac{6+8+10+12+14}{5} &&\\
\overline{x} &= \dfrac{50}{5} &&\\
\overline{x} &= \res{10} &&
\end{calculs}
\begin{center}
\renewcommand{\arraystretch}{1.5}
\begin{tabular}{|l|*{5}{c|}}
\hline
Pages $x$ & $6$ & $8$ & $10$ & $12$ & $14$\\\hline
Écart $x-10$ & $-4$ & $-2$ & $0$ & $2$ & $4$\\\hline
Carré de l'écart & $16$ & $4$ & $0$ & $4$ & $16$\\\hline
\end{tabular}
\end{center}
\begin{calculs}
\sigma &= \sqrt{\dfrac{16+4+0+4+16}{5}} &&\justif{moyenne des carrés, puis racine}\\
\sigma &= \sqrt{8} &&\\
\sigma &\approx \res{2{,}83} &&\justif{pages}
\end{calculs}
\cas{b)}
\explique{Application \texttt{Statistiques} : $\overline{x}=\res{10}$ et
$\sigma\approx\res{6{,}32}$.}{les cinq valeurs dans \texttt{V1}, $1$ dans chaque case de
\texttt{N1}}
\cas{c)}
\explique{Léo et Zoé lisent en moyenne $10$ pages par jour, mais l'écart type de Zoé est
plus du double : \res{ses lectures sont bien plus irrégulières}.}{même moyenne,
dispersions différentes}
\end{exo}

\begin{exo}[Lancers de trois pièces]
\cas{a)}
\begin{calculs}
\overline{x} &= \dfrac{3\times 0+9\times 1+9\times 2+3\times 3}{24} &&\\
\overline{x} &= \dfrac{36}{24} &&\\
\overline{x} &= \res{1{,}5} &&
\end{calculs}
\cas{b)}
\explique{\res{$\sigma=\sqrt{\dfrac{3\times(0-1{,}5)^{2}+9\times(1-1{,}5)^{2}+9\times(2-1{,}5)^{2}
+3\times(3-1{,}5)^{2}}{24}}$}}{chaque carré d'écart est multiplié par son effectif}
\cas{c)}
\begin{calculs}
\sigma &= \sqrt{\dfrac{3\times 2{,}25+9\times 0{,}25+9\times 0{,}25+3\times 2{,}25}{24}}
  &&\justif{$(-1{,}5)^{2}=2{,}25$ ; $(-0{,}5)^{2}=0{,}25$}\\
\sigma &= \sqrt{\dfrac{6{,}75+2{,}25+2{,}25+6{,}75}{24}} &&\\
\sigma &= \sqrt{\dfrac{18}{24}} &&\\
\sigma &= \sqrt{0{,}75} &&\\
\sigma &\approx \res{0{,}87} &&\justif{la calculatrice affiche $\sigma\approx 0{,}866$}
\end{calculs}
\end{exo}

\begin{exo}[Choisir avec les indicateurs]
\cas{a)}
\explique{\res{La machine $A$ est la plus régulière} : les deux moyennes sont proches de
$250$ g, mais son écart type est bien plus petit.}{à moyennes voisines, on compare les
écarts types}
\cas{b)}
\explique{En \res{seconde E} : même moyenne, mais un écart type plus grand ($5>2$).}{plus
l'écart type est grand, plus les valeurs sont dispersées}
\cas{c)}
\explique{Médiane : \res{au moins la moitié des salariés gagnent au plus $2~100$ €}, et au
moins la moitié au moins $2~100$ €.}{la médiane partage la série en deux}
\explique{Écart interquartile : au moins la moitié des salaires sont compris entre $Q_1$ et
$Q_3$, dans un intervalle de largeur $300$ €.}{$\intff{Q_1}{Q_3}$ contient au moins la
moitié des valeurs}
\end{exo}

\begin{exo}[Sélection au tir à l'arc]
\cas{a)}
\begin{calculs}
\overline{x}_X &= \dfrac{1\times 0+3\times 4+6\times 6+7\times 8+3\times 10}{20} &&\\
\overline{x}_X &= \dfrac{134}{20} &&\\
\overline{x}_X &= \res{6{,}7} &&
\end{calculs}
\begin{calculs}
\overline{x}_Y &= \dfrac{3\times 0+3\times 4+4\times 6+4\times 8+6\times 10}{20} &&\\
\overline{x}_Y &= \dfrac{128}{20} &&\\
\overline{x}_Y &= \res{6{,}4} &&
\end{calculs}
\cas{b)}
\explique{$N=20$ : rangs $10$ et $11$. Pour $X$, effectifs cumulés $1$ ; $4$ ; $10$ ;
$17$ ; $20$ : la $10^{\text{e}}$ valeur est $6$, la $11^{\text{e}}$ est $8$, donc
$Me=\res{7}$.}{partie III}
\explique{Pour $Y$, effectifs cumulés $3$ ; $6$ ; $10$ ; $14$ ; $20$ : de même
$Me=\dfrac{6+8}{2}=\res{7}$.}{les deux médianes sont égales}
\cas{c)}
\explique{Calculatrice : $\sigma_X\approx\res{2{,}39}$ et $\sigma_Y\approx\res{3{,}38}$.}{
valeurs dans \texttt{V1}, nombres de flèches dans \texttt{N1}}
\cas{d)}
\explique{\res{L'archer $X$ est le plus régulier} : son écart type est plus petit, avec
une moyenne un peu meilleure et la même médiane. On choisit $X$.}{la régularité se
mesure par l'écart type}
\end{exo}

\partie{Diagrammes}

\begin{exo}[Matière préférée]
\cas{a)}
\explique{Les deux « ne se prononce pas » sont retirés : $N=6+5+4+7+8=30$.}{effectif
total}
\begin{calculs}
\text{angle} &= \dfrac{6}{30}\times 360 &&\justif{maths : fréquence $\times 360^{\circ}$}\\
\text{angle} &= 0{,}2\times 360 &&\justif{fréquence $20\,\%$}\\
\text{angle} &= \res{72^{\circ}} &&
\end{calculs}
\explique{Chaque élève compte pour $\dfrac{360}{30}=12^{\circ}$ : langues
$60^{\circ}$ ($16{,}7\,\%$), histoire $48^{\circ}$ ($13{,}3\,\%$), sciences
$84^{\circ}$ ($23{,}3\,\%$), EPS $96^{\circ}$ ($26{,}7\,\%$).}{vérification :
$72+60+48+84+96=360$}
\cas{b)}
\begin{center}
\begin{tikzpicture}
  \fill[coulCourbe!20] (0,0) -- (90:2) arc (90:162:2) -- cycle;
  \fill[coulCourbeC!20] (0,0) -- (162:2) arc (162:222:2) -- cycle;
  \fill[coulCourbeB!20] (0,0) -- (222:2) arc (222:270:2) -- cycle;
  \fill[coulCourbe!8] (0,0) -- (270:2) arc (270:354:2) -- cycle;
  \fill[coulCourbeC!8] (0,0) -- (354:2) arc (354:450:2) -- cycle;
  \draw[secteur] (0,0) circle (2);
  \foreach \a in {90,162,222,270,354} \draw[secteur] (0,0) -- (\a:2);
  \node[font=\footnotesize] at (126:1.3) {maths};
  \node[font=\footnotesize] at (192:1.3) {langues};
  \node[font=\footnotesize] at (246:1.3) {histoire};
  \node[font=\footnotesize] at (312:1.25) {sciences};
  \node[font=\footnotesize] at (42:1.2) {EPS};
\end{tikzpicture}
\end{center}
\cas{c)}
\explique{Un \res{diagramme en barres}. Pas d'histogramme : le caractère est qualitatif,
ses valeurs ne sont pas des nombres regroupés en classes.}{définition 10}
\end{exo}

\begin{exo}[Lire un diagramme circulaire]
\cas{a)}
\begin{calculs}
\text{football} &= \dfrac{120}{360}\times 240 &&\justif{l'angle est proportionnel à l'effectif}\\
\text{football} &= \res{80} &&
\end{calculs}
\explique{De même : basket $\dfrac{60}{360}\times 240=40$ ; tennis $30$ ; handball
$60$ ; escalade $30$.}{vérification : $80+40+30+60+30=240$}
\cas{b)}
\explique{$\dfrac{90}{360}=0{,}25$ : \res{$25\,\%$} des adhérents pratiquent le
handball.}{un quart du disque}
\end{exo}

\begin{exo}[Histogramme et polygone]
\cas{a)}
\par\noindent
\begin{minipage}[t]{0.55\linewidth}\raggedright
Axe horizontal : $1$ carreau pour $7{,}5$ min ; axe vertical : $1$ carreau
pour $2$ salariés. Rectangles accolés, de largeur $15$ min, de hauteur égale à
l'effectif.
\end{minipage}\hfill
\begin{minipage}[t]{0.42\linewidth}\vspace{0pt}\centering
\begin{tikzpicture}[x=0.5cm,y=0.3cm]
  \quadrillage{0}{0}{9}{16}
  \foreach \a/\h in {0/9,2/15,4/12,6/4} \draw[barre] (\a,0) rectangle ({\a+2},\h);
  \draw[-{Stealth[length=2mm]},thick] (0,0)--(9.5,0) node[above left,font=\footnotesize]{min};
  \draw[-{Stealth[length=2mm]},thick] (0,0)--(0,16.6) node[left,font=\footnotesize]{effectif};
  \foreach \k/\l in {0/0,2/15,4/30,6/45,8/60}
    \draw (\k,0.25)--(\k,-0.25) node[below,font=\scriptsize]{$\l$};
  \foreach \k/\l in {5/10,10/20,15/30}
    \draw (0.12,\k)--(-0.12,\k) node[left,font=\scriptsize]{$\l$};
\end{tikzpicture}
\end{minipage}\par

\cas{b)}
\explique{Effectifs cumulés $18$ ; $48$ ; $72$ ; $80$, divisés par $80$, donnent $22{,}5\,\%$ ;
$60\,\%$ ; $90\,\%$ ; \res{$100\,\%$}.}{fréquences cumulées croissantes}
\cas{c)}
\par\noindent
\begin{minipage}[t]{0.55\linewidth}\raggedright
On place $(0\,;0)$, $(15\,;22{,}5)$, $(30\,;60)$, $(45\,;90)$ et $(60\,;100)$ :
chaque fréquence cumulée à la fin de sa classe, puis on relie par des segments.

\smallskip
\textbf{\color{coulExo}d)} L'horizontale d'ordonnée $50\,\%$ coupe le polygone au-dessus
de $Me\approx\res{26}$ min.
\end{minipage}\hfill
\begin{minipage}[t]{0.42\linewidth}\vspace{0pt}\centering
\begin{tikzpicture}[x=0.5cm,y=0.4cm]
  \quadrillage{0}{0}{9}{10}
  \draw[-{Stealth[length=2mm]},thick] (0,0)--(9.5,0) node[above left,font=\footnotesize]{min};
  \draw[-{Stealth[length=2mm]},thick] (0,0)--(0,10.6) node[right,font=\footnotesize]{\%};
  \foreach \k/\l in {0/0,2/15,4/30,6/45,8/60}
    \draw (\k,0.15)--(\k,-0.15) node[below,font=\scriptsize]{$\l$};
  \foreach \k/\l in {2/20,5/50,8/80,10/100}
    \draw (0.12,\k)--(-0.12,\k) node[left,font=\scriptsize]{$\l$};
  \draw[red,dashed,thick] (0,5)--(3.47,5)--(3.47,0);
  \draw[coulCourbe,line width=1.1pt] (0,0)--(2,2.25)--(4,6)--(6,9)--(8,10);
  \foreach \p in {(0,0),(2,2.25),(4,6),(6,9),(8,10)}
    \path[coulCourbe] \p node[croix]{};
\end{tikzpicture}
\end{minipage}\par

\cas{e)}
\explique{$\dfrac{80}{2}=40$ et $18<40\le 48$ : classe médiane $\intfo{15}{30}$. Il faut
$40-18=22$ salariés de cette classe, sur ses $30$.}{méthode 6}
\begin{calculs}
Me &\approx 15+\dfrac{22}{30}\times 15 &&\\
Me &\approx 15+11 &&\\
Me &\approx \res{26} &&\justif{on retrouve la lecture graphique}
\end{calculs}
\end{exo}

\begin{exo}[Lire un histogramme]
\cas{a)}
\explique{On lit la hauteur de chaque rectangle : $5$ ; $12$ ; $9$ ; $6$ ; $3$.}{$1$
carreau pour $1$ élève}
\begin{center}
\renewcommand{\arraystretch}{1.5}
\begin{tabular}{|l|*{5}{c|}}
\hline
Temps de sport (h) & $\intfo{0}{2}$ & $\intfo{2}{4}$ & $\intfo{4}{6}$ & $\intfo{6}{8}$
  & $\intfo{8}{10}$\\\hline
Effectif & $5$ & $12$ & $9$ & $6$ & $3$\\\hline
\end{tabular}
\end{center}
\explique{Le tableau reprend la lecture.}{\res{$5$ ; $12$ ; $9$ ; $6$ ; $3$}}
\cas{b)}
\explique{$5+12+9+6+3=\res{35}$ élèves.}{somme des effectifs}
\cas{c)}
\explique{Au moins $6$ h : $6+3=9$ élèves, et $\dfrac{9}{35}\approx 0{,}257$, soit
\res{$25{,}7\,\%$}.}{classes $\intfo{6}{8}$ et $\intfo{8}{10}$}
\cas{d)}
\begin{calculs}
\overline{x} &\approx \dfrac{5\times 1+12\times 3+9\times 5+6\times 7+3\times 9}{35}
  &&\justif{centres des classes}\\
\overline{x} &\approx \dfrac{5+36+45+42+27}{35} &&\\
\overline{x} &\approx \dfrac{155}{35} &&\\
\overline{x} &\approx \res{4{,}43} &&\justif{heures par semaine}
\end{calculs}
\end{exo}

\partie{Croisement de deux variables qualitatives}

\begin{exo}[Petit-déjeuner]
\cas{a)}
\begin{calculs}
\text{oui, première} &= 0{,}65\times 260 &&\justif{$65\,\%$ de la colonne « première »}\\
\text{oui, première} &= \res{169} &&
\end{calculs}
\explique{Non, première : $260-169=91$ ; puis les totaux des lignes.}{chaque colonne a
pour somme son total}
\begin{center}
\renewcommand{\arraystretch}{1.5}
\begin{tabular}{|l|c|c|c|c|}
\hline
 & seconde & première & terminale & Total\\\hline
Oui & $210$ & $169$ & $156$ & $535$\\\hline
Non & $90$ & $91$ & $84$ & $265$\\\hline
Total & $300$ & $260$ & $240$ & $800$\\\hline
\end{tabular}
\end{center}
\cas{b)}
\begin{calculs}
f &= \dfrac{535}{800} &&\justif{total de la ligne « oui » sur l'effectif total}\\
f &\approx \res{0{,}669} &&\justif{soit environ $66{,}9\,\%$}
\end{calculs}
\cas{c)}
\begin{calculs}
f &= \dfrac{156}{240} &&\justif{on ne regarde que la colonne « terminale »}\\
f &= \res{0{,}65} &&\justif{soit $65\,\%$}
\end{calculs}
\cas{d)}
\begin{calculs}
f &= \dfrac{84}{265} &&\justif{on ne regarde que la ligne « non »}\\
f &\approx \res{0{,}317} &&\justif{soit environ $31{,}7\,\%$}
\end{calculs}
\explique{Environ $32\,\%$ des élèves qui n'ont pas pris de petit-déjeuner sont en
terminale.}{population de référence : ceux qui n'ont pas déjeuné}
\end{exo}

\begin{exo}[Lire les phrases des médias]
\cas{a)}
\explique{(1) Fréquence \res{conditionnelle} ; population de référence : les actifs qui
habitent une ville-centre.}{« des actifs qui habitent\ldots{} » : une partie des actifs}
\explique{(2) Fréquence conditionnelle ; population de référence : les actifs des
communes hors de l'attraction des villes.}{« chez les actifs de\ldots{} »}
\explique{(3) Fréquence \res{marginale} ; population de référence : tous les
actifs.}{aucune condition}
\cas{b)}
\explique{La phrase (1) parle de la part des cyclistes parmi les habitants des
villes-centres ; l'article, de la part des habitants des villes-centres \res{parmi les
cyclistes} : la population de référence a changé.}{« $x\,\%$ des A sont B » n'est pas
« $x\,\%$ des B sont A »}
\end{exo}

\begin{exo}[Le vélo et la distance]
\cas{a)}
\begin{calculs}
\text{$2$ km ou moins} &= 0{,}05\times 3~700 &&\justif{$5\,\%$ de $3~700$ milliers}\\
\text{$2$ km ou moins} &= 185 &&
\end{calculs}
\explique{De $2$ à $5$ km : $0{,}051\times 4~500=229{,}5\approx 230$ ; plus de $5$ km :
$0{,}009\times 16~400=147{,}6\approx 148$.}{arrondis à l'unité}
\explique{Total : $185+230+148=\res{563}$ milliers de cyclistes.}{somme des trois}
\cas{b)}
\begin{calculs}
f &= \dfrac{563}{24~600} &&\justif{cyclistes sur l'ensemble des actifs}\\
f &\approx 0{,}0229 &&\\
f &\approx \res{2{,}3\,\%} &&
\end{calculs}
\cas{c)}
\begin{calculs}
f &= \dfrac{185+230}{563} &&\justif{parmi les cyclistes : on divise par $563$}\\
f &= \dfrac{415}{563} &&\\
f &\approx 0{,}737 &&
\end{calculs}
\explique{\res{Près des trois quarts ($73{,}7\,\%$) des cyclistes habitent à $5$ km ou
moins de leur travail.}}{alors que seulement $5\,\%$ des actifs proches viennent à vélo}
\end{exo}

\begin{exo}[Satisfaction sur une ligne de bus]
\cas{a)}
\explique{Occasionnels : $400-250=150$. Réguliers satisfaits : $0{,}7\times 250=175$ ;
occasionnels satisfaits : $0{,}8\times 150=120$ ; le reste par soustraction.}{pourcentages
de chaque ligne}
\begin{center}
\renewcommand{\arraystretch}{1.5}
\begin{tabular}{|l|c|c|c|}
\hline
 & Satisfait & Non satisfait & Total\\\hline
Régulier & $175$ & $75$ & $250$\\\hline
Occasionnel & $120$ & $30$ & $150$\\\hline
Total & $295$ & $105$ & $400$\\\hline
\end{tabular}
\end{center}
\explique{Le tableau donne $\res{295}$ usagers satisfaits.}{total de la colonne}
\cas{b)}
\begin{calculs}
f &= \dfrac{295}{400} &&\justif{total de la colonne sur l'effectif total}\\
f &= \res{0{,}7375} &&\justif{soit $73{,}75\,\%$}
\end{calculs}
\cas{c)}
\begin{calculs}
f &= \dfrac{75}{105} &&\justif{on ne regarde que la colonne « non satisfait »}\\
f &\approx \res{0{,}714} &&\justif{soit environ $71{,}4\,\%$}
\end{calculs}
\cas{d)}
\explique{\res{Non} : $70\,\%$ est la part des satisfaits parmi les réguliers seulement ;
parmi tous les usagers, elle vaut $73{,}75\,\%$.}{le journal confond fréquence
conditionnelle et fréquence marginale}
\cas{e)}
\explique{Usager choisi au hasard : $P(R\cap S)=\dfrac{175}{400}=\res{0{,}4375}$.}{la case
« régulier, satisfait »}
\explique{$P_S\big(\overline{R}\big)=\dfrac{120}{295}=\dfrac{24}{59}\approx\res{0{,}41}$.}{
parmi les satisfaits : la colonne « satisfait » (chapitre 11)}
\end{exo}

\end{document}
